% SEGMENT OLP-0344-S01
% Part: lambda-calculus
% Chapter: introduction
% Section: syntax

\documentclass[../../../include/open-logic-section]{subfiles}

\begin{document}

\olfileid{lam}{int}{syn}
\olsection{லாம்டா கலனத்தின் அமைப்புவிதி}

$x$, $y$, $z$,~\dots என்ற மாறிகளின் தொடர்முறையுடனும், $a$, $b$, $c$,~\dots
என்ற சில மாறிலிக் குறிகளுடனும் தொடங்குகிறோம். சொற்களின் கணம் பின்வருமாறு
தொகுத்தறிதல் முறையில் வரையறுக்கப்படுகிறது:
\begin{enumerate}
\item ஒவ்வொரு மாறியும் ஒரு சொல்.
\item ஒவ்வொரு மாறிலியும் ஒரு சொல்.
\item $M$, $N$ ஆகியவை சொற்கள் எனில், $(MN)$ உம் ஒரு சொல்.
\item $M$ ஒரு சொல்லாகவும் $x$ ஒரு மாறியாகவும் இருந்தால்,
  $(\lambd[x][M])$ உம் ஒரு சொல்.
\end{enumerate}
$(MN)$ என்ற வடிவிலுள்ள சொற்கள் \emph{செயல்படுத்தல்கள்} என்றும்,
$(\lambd[x][M])$ என்ற வடிவிலுள்ளவை \emph{அருவமாக்கங்கள்} என்றும் அழைக்கப்படுகின்றன.

மாறிலிகள் எவையும் இல்லாத முறைமை \emph{தூய} லாம்டா கலனம் எனப்படுகிறது. நாம்
பெரும்பாலும் தூய $\lambd$-கலனத்திலேயே பணிபுரிவோம்; எனவே சிற்றெழுத்துகள் அனைத்தும்
மாறிகளைக் குறிக்கும். $\lambd$-கலனத்தின் சொற்களைக் குறிக்கப் பேரெழுத்துகளை
($M$, $N$ முதலியவற்றை) பயன்படுத்துகிறோம்.

பின்வரும் சில குறிமான மரபுகளைப் பின்பற்றுவோம்:

\begin{conv}
\begin{enumerate}
\item அடைப்புக்குறிகள் விடப்பட்டால், செயல்படுத்தல் இடமிருந்து வலமாக நிகழும்.
  எடுத்துக்காட்டாக, $M$, $N$, $P$, $Q$ ஆகியவை சொற்கள் எனில், $MNPQ$ என்பது
  $(((MN)P)Q)$ என்பதன் சுருக்கம்.
% SOURCE CORRECTION TA-LAM-001: read the frozen prose typo "From example" as "For example."
\item மீண்டும், அடைப்புக்குறிகள் விடப்பட்டால், லாம்டா அருவமாக்கத்திற்குச்
  சாத்தியமான மிகப்பரந்த வீச்சு வழங்கப்படும். எடுத்துக்காட்டாக,
  $\lambd[x][MNP]$ என்பது $(\lambd[x][((MN)P)])$ என்று வாசிக்கப்படும்.
\item பல மாறிகளை அருவமாக்க ஒரே லாம்டாவைப் பயன்படுத்தலாம். எடுத்துக்காட்டாக,
  $\lambd[xyz][M]$ என்பது
  $\lambd[x][\lambd[y][\lambd[z][M]]]$ என்பதன் சுருக்கம்.
\end{enumerate}
\end{conv}

எடுத்துக்காட்டாக,
\[
\lambd[xy][xxyx \lambd[z][xz]]
\]
என்பது
\[
\lambd[x][\lambd[y][((((xx)y)x)(\lambd[z][(xz)]))]].
\]
என்பதன் சுருக்கம். இந்த மரபுகளை நீங்கள் மனப்பாடம் செய்ய வேண்டும். முதலில் அவை
உங்களை மிகவும் குழப்பும்; ஆனால் அவற்றுக்குப் பழகிவிடுவீர்கள், சிறிது காலத்திற்குப்
பிறகு அடைப்புக்குறிகளின் புதரைக் கையாள்வதைவிட அவை குறைவாகவே குழப்பும்.

கட்டுண்ட மாறிகளின் பெயர்களில் மட்டுமே வேறுபடும் இரு சொற்கள்
$\alpha$-சமானமானவை எனப்படுகின்றன; எடுத்துக்காட்டாக,
$\lambd[x][x]$, $\lambd[y][y]$. இவற்றை ``ஒரே'' சொல்லாகக் கருதுவது வசதியானது;
வேறு சொற்களில், $M$, $N$ இரண்டும் ஒரேவை என்று நாம் கூறும்போது, ``கட்டுண்ட
மாறிகளின் மறுபெயரிடல்கள் வரையிலான ஒரேவை'' என்றும் பொருள்கொள்கிறோம்.
% SOURCE CORRECTION TA-LAM-002: an occurrence is bound only by an enclosing lambda for the same variable; the nearest matching binder applies.
ஒரு மாறியின் நிகழ்வு, அதே மாறியைப் பிணைக்கும் ஓர் வெளிச்சூழ்
$\lambd$-அருவமாக்கத்தின் வீச்சுக்குள் இருந்தால் அது ``கட்டுண்டது''; அத்தகைய
பொருந்தும் பிணைப்பிகள் ஒன்றுக்கு மேல் இருந்தால், மிக அருகிலுள்ள பிணைப்பியே அந்த
நிகழ்வைப் பிணைக்கிறது. பொருந்தும் வெளிச்சூழ்ப் பிணைப்பி இல்லாத நிகழ்வு
``கட்டற்றது'' என்று அழைக்கப்படுகிறது. முந்தைய எடுத்துக்காட்டில் கட்டற்ற
மாறிநிகழ்வுகள் இல்லை; ஆனால்
\[
(\lambd[z][yz])x
\]
என்பதில் $y$, $x$ ஆகியவை கட்டற்றவை; $z$ கட்டுண்டது.

\end{document}
